%% manus/main.tex — STAGE manuscript entry point (one paper = one repo).
%%
%% Build: bash execs/run.sh          latexmk, out-of-tree; the PDF lands in
%%                                   wkdrs/builds/, never next to the sources.
%% Lint:  bash execs/scpts/lint.sh   deterministic checks; any todo marker in
%%                                   manus/ fails it — that is the point.
%%
%% Venue class: `article` is a standalone-compiling placeholder. To adopt a
%% venue template, drop its class/style files into manus/stys/ and swap the
%% line below, e.g. \documentclass{stys/cvpr}. run.sh compiles with manus/ as
%% the working directory and puts manus/stys/ on TEXINPUTS.
\documentclass{article}

%% STAGE base macros: the todo marker plus graphicx/booktabs/hyperref/xcolor.
\usepackage{stys/stage}

\title{Untitled STAGE Manuscript}
%% Stays "Anonymous" while ANON=true in .env; lint.sh hunts identity leaks.
\author{Anonymous}
\date{}

\begin{document}
\maketitle

\begin{abstract}
%% Placeholder — /stage-stry-coach shapes the story behind this paragraph;
%% /stage-plan-outliner may move the abstract into secs/0_abstract.tex.
This manuscript was scaffolded by STAGE, the academic-writing companion to
STAR. Every number it will ever state either traces to fingerprinted evidence
under \texttt{mates/} or stays a visible todo marker; the placeholder below
demonstrates the second state, and the skeleton is expected to fail lint until
it is gone.
\end{abstract}

%% Placeholder body — superseded by the secs/ inputs below as they come alive.
\section{Introduction}
\todo{No content yet: run /stage-stry-coach, then /stage-plan-outliner.}

%% Section files — /stage-plan-outliner creates manus/secs/<n>_<slug>.tex
%% skeletons and uncomments one \input line per created file, in the order
%% notes/outline.md fixes. Example shape:
% \input{secs/1_intro}
% \input{secs/2_related}
% \input{secs/3_method}
% \input{secs/4_experiments}
% \input{secs/5_conclusion}

%% Bibliography — uncomment once manus/bibs/reference.bib exists
%% (execs/scpts/import.sh seeds it from a paired STAR repo, or
%% /stage-refs-curator starts one from scratch).
% \bibliographystyle{plain}
% \bibliography{bibs/reference}

\end{document}
